\begin{abstract}
Memory is essential for agents to retain past experiences and use them in future interactions. Recent studies show that, as models become more capable, simple retrieval-augmented generation (RAG) can answer most text-memory questions correctly and outperform most memory systems. However, this level of performance does not consistently carry over to multimodal memory. One reason is that visual evidence often spans several related images, some of which contain crucial information but have low embedding similarity to the question. Similarity-based retrieval can overlook these images, leaving the answering model without the evidence needed to connect relevant facts and derive the answer. Moreover, each image can contain rich information, so missing even one relevant image can leave a substantial gap in the evidence. These challenges make reliable multimodal memory difficult to achieve through similarity-based retrieval alone. To address this problem, we propose Visibility-Gated Set Reading (\vgs{}), which combines global evidence navigation with visibility-guided reading. Global navigation reads a compact history index to select related rounds together as an evidence set. The visibility gate widens evidence access for histories whose content is poorly represented in the index, using a reading path suited to the answering model. This helps \vgs{} recover related images together and retain the evidence needed for reasoning. Experiments on four multimodal memory benchmarks show that \vgs{} improves a lightweight reader over RAG on all four, by 4 to 13 points, and a strong reader on long histories, where most image evidence lies outside the retrieved rounds.
\end{abstract}
